\documentclass[11pt]{article}
\usepackage[margin=1in]{geometry}
\usepackage{amsmath,amssymb,amsthm,mathtools}
\usepackage{booktabs,tabularx,array}
\usepackage{microtype}
\usepackage{enumitem}
\usepackage{xurl}
\usepackage[colorlinks=true,linkcolor=blue,citecolor=blue,urlcolor=blue]{hyperref}
\providecommand{\doi}[1]{\href{https://doi.org/#1}{doi:#1}}

\newcommand{\R}{R(3,3,3,3)}
\newcommand{\Nc}{N_c}
\newtheorem{theorem}{Theorem}[section]
\newtheorem{lemma}[theorem]{Lemma}
\newtheorem{proposition}[theorem]{Proposition}
\newtheorem{corollary}[theorem]{Corollary}
\newtheorem{computation}[theorem]{Finite computation}
\theoremstyle{definition}
\newtheorem{definition}[theorem]{Definition}
\theoremstyle{remark}
\newtheorem{remark}[theorem]{Remark}
\usepackage[nameinlink]{cleveref}
\crefname{computation}{finite computation}{finite computations}
\Crefname{computation}{Finite computation}{Finite computations}

\title{A certified common-critical obstruction in\\
four-colour triangle-free graphs\\
\large With an audit of a 113,179-case reduction at 58}
\author{The r3333-58 project\\[-1pt]
\small Working draft; author list and affiliations to be supplied}
\date{25 November 2026\\[-1pt]
\small Mathematical source snapshot: repository revision \texttt{35140c0}}

\begin{document}
\maketitle

\begin{abstract}
\noindent
Call an edge-colouring of a complete graph \emph{good} if it has no
monochromatic triangle, and call a colour \emph{critical} at a vertex when its
degree there is at least $15$.  We give a self-contained reduction, backed by
496 replayed DRAT refutations, showing that two distinct colours cannot both
be critical at every vertex of a nonempty good four-colouring.  The certificate
uses catalogues on $15$ and $16$ vertices and two encodings of a
25-vertex obstruction.  It does not depend on the semantic bridge for the
larger $K_{58}$ SAT cover.

For a good four-colouring of $K_{58}$, assume that criticality is reciprocal
along every edge.  Connectivity then makes criticality constant in each colour,
and the obstruction leaves at most the two sorted degree profiles
$(15,14,14,14)$ and $(16,14,14,13)$.  This is a restriction of the reciprocal
case, not an elimination of unrestricted colourings.  A further counting
argument gives either zero or at least $15$ unreciprocated tickets, and at
least $30$ under vertexwise ticket balance.

The second part records the associated formalisation and evidence audit.  Lean
proves the signed-unit cover completeness statement for 113,179 frozen cases.
The complete encoding bridge and the computational premise remain unfinished.
In the 25 November evidence snapshot, the owner reports 98,142 replayed case
certificates; 12,893 rows are OPEN and 2,144 require legacy evidence.  The
proof bytes for those 98,142 cases were not included in that release.  Thus no
new bound on $\R$ is claimed here.
\end{abstract}

\section{Introduction}\label{sec:intro}

The multicolour Ramsey number $\R$ is the least $n$ for which every
four-colouring of the edges of $K_n$ contains a monochromatic triangle.  The
published bounds are $51\leq\R\leq62$, due respectively to Chung and to
Fettes, Kramer and Radziszowski
\cite{Chung1973,FKR2004,Radziszowski2026}.  A recent arXiv preprint by
Tatarevic claims the improvement $\R\leq60$ using a global ticket count,
small-neighbourhood classifications and exhaustive SAT computations, with the
mathematical reductions formalised in Lean \cite{Tatarevic2026}.  We cite that
claim as a preprint and do not treat it as one of the results established here.

This paper has two connected purposes.  The first is a general obstruction for
good four-colourings.  Write $d_c(v)$ for the number of colour-$c$ edges at
$v$.  We call $c$ \emph{critical} at $v$ if $d_c(v)\geq15$.

\begin{theorem}[Common-critical obstruction]\label{thm:common}
In every nonempty good four-colouring of a finite complete graph, and for any
two distinct colours $a$ and $b$, there is a vertex $x$ such that
$d_a(x)\leq14$ or $d_b(x)\leq14$.  Equivalently, at most one colour can be
critical at every vertex.
\end{theorem}

\Cref{thm:common} is finite and certified, rather than an inference from a
solver timeout or a bounded search.  Its reduction uses a 25-vertex local
obstruction.  Two catalogue refutations, 247 refutations of a slot encoding
and 247 refutations of a word encoding establish the finite residue.  The
complete certificate was replayed during preparation of this draft.
\Cref{sec:certificate} describes what the encodings prove and, equally
important, why a physical colouring would extend to their auxiliary variables.

The second purpose is to make the evidence boundary of the $K_{58}$ campaign
legible.  A signed-unit cover theorem, a DIMACS file, a solver exit code and a
replayed refutation are different objects.  They must be joined in the right
direction before a Ramsey bound follows.  The repository contains a Lean
formalisation of cover completeness at the signed-unit level, but not yet the
full bridge from a physical colouring to every auxiliary-variable clause in
the emitted formulas.  It also does not contain a complete set of case
refutations.  We therefore state explicitly:
\begin{center}
\emph{This paper does not prove $\R\leq58$.}
\end{center}

\subsection{Consequences for reciprocal colourings}\label{sec:consequences}

A good four-colouring of $K_{58}$ is \emph{reciprocal} when, for every edge
$uv$, the colour of $uv$ is critical at $u$ exactly when it is critical at
$v$.  In the ticket terminology of \cite{Tatarevic2026}, this says that every
ticket is returned.

\begin{corollary}[Reciprocal profiles at $58$]\label[corollary]{cor:profiles}
If a good four-colouring of $K_{58}$ is reciprocal, every vertex has the same
sole common critical colour, and its sorted degree profile is either
\[
 (15,14,14,14)\quad\text{or}\quad(16,14,14,13).
\]
No assertion of existence for either profile is made, and nothing is concluded
about a nonreciprocal colouring.
\end{corollary}

Thus, conditional on reciprocity, nine of the eleven degree profiles in the
frozen $K_{58}$ analysis are impossible.  This corollary supplies no DRAT
refutation of the original, unrestricted CNFs for their rows.  The word
``conditional'' cannot be dropped: an unreciprocated ticket can allow
criticality to change along a colour component.

The same local obstruction gives information about that nonreciprocity.
Define the high and low sets for a colour $c$ by
\[
 H_c=\{v:d_c(v)\geq15\},\qquad X_c=V\setminus H_c.
\]
A colour-$c$ edge from $H_c$ to $X_c$ carries one unreciprocated ticket.

\begin{corollary}[Ticket gap]\label[corollary]{cor:ticket-gap}
Every good four-colouring of $K_{58}$ has either zero unreciprocated tickets or
at least $15$.  If two distinct colours are critical somewhere, the second
alternative holds.  If additionally every vertex has equal total incoming and
outgoing ticket counts, there are at least $30$ unreciprocated tickets.
\end{corollary}

As a separate observation, \Cref{thm:common} also shortens the reciprocal
part of the perfectly balanced argument in \cite{Tatarevic2026}: on $K_{60}$
every feasible degree profile has at least two critical colours, so a global
reciprocity hypothesis would make two colours critical everywhere.  This does
not reproduce or replace the nonreciprocal computations in that proof.

\subsection{Organisation}

\Cref{sec:local,sec:certificate} prove \Cref{thm:common} modulo the stated
finite computations and explain their semantics.  \Cref{sec:reciprocity}
proves the profile restriction, and \Cref{sec:ticket-gap} proves the ticket
gap.  \Cref{sec:cover,sec:formal,sec:status} describe the frozen reduction,
the Lean boundary and the current computational status.  \Cref{sec:audit}
records the verification methodology and the claims that remain open.

\section{The 25-vertex obstruction}\label{sec:local}

We use $R(3,3)=6$ and $R(3,3,3)=17$ \cite{GreenwoodGleason1955}.  A colour
neighbourhood in a good colouring omits its own colour, since an internal edge
of that colour would close a monochromatic triangle.  It therefore carries a
good three-colouring and has at most $16$ vertices.  In particular, inside a
colour neighbourhood of size at least $6$, every vertex has at most five
neighbours in any one of the remaining colours.

\begin{lemma}[Extendible catalogues]\label[lemma]{lem:extend}
Every good three-colouring of $K_{15}$ extends, as a coloured graph, to a good
three-colouring of $K_{16}$.  In a good three-colouring of $K_{16}$, every
vertex has degree five in each colour.  Consequently, after deleting the added
vertex, ten vertices of the original $K_{15}$ retain degree five in any chosen
colour.
\end{lemma}

The degree assertion follows immediately: degrees are at most five and sum to
$15$.  Extendibility is a finite catalogue fact, checked by
\Cref{comp:catalogues} below.  It is also consistent with the known two-template
classification of good three-colourings at orders $15$ and $16$
\cite{KalbfleischStanton1968,Heinrich1977}, but the present certificate does not
need to assume an external completeness argument: its two catalogue CNFs close
that finite step directly.

\begin{proposition}[Local obstruction]\label[proposition]{prop:25}
Let $a,b$ be distinct colours and let $L,R$ be 15-vertex sets in a good
four-colouring such that $|L\cap R|=5$.  Suppose that $L$ has no colour-$a$
edge and $R$ has no colour-$b$ edge.  Then the following two assertions cannot
both hold:
\begin{enumerate}[label=\textup{(\roman*)},nosep]
 \item every vertex of $L$ has ambient colour-$b$ degree at least $15$;
 \item every vertex of $R$ has ambient colour-$a$ degree at least $15$.
\end{enumerate}
Here ``ambient'' refers to the complete host graph; only consequences on the
25 named vertices are used in the finite encoding.
\end{proposition}

\Cref{prop:25} is exactly the content certified by the slot and word
refutations in \Cref{comp:orbits}.  We first use it to prove the global
statement.

\begin{proof}[Proof of \Cref{thm:common}]
Suppose, for contradiction, that $a$ and $b$ are both critical at every
vertex.  Choose $r$ and select $15$ colour-$a$ neighbours of $r$ to form $L$.
Then $L$ omits colour $a$ and hence carries a good three-colouring.

By \Cref{lem:extend}, add one vertex to obtain a good three-coloured $K_{16}$.
Each vertex of this extension has degree five in each of its three colours.
Deleting the added vertex leaves ten vertices of $L$ with exactly five
colour-$b$ neighbours inside $L$.  Choose such a vertex $y$.

Since $b$ is critical at $y$ in the ambient graph, choose $15$ of its ambient
colour-$b$ neighbours, including its five colour-$b$ neighbours in $L$, and
call the resulting set $R$.  Then $R$ omits $b$ and $L\cap R$ consists exactly
of those five vertices.  Moreover, every vertex of $L$ is $b$-critical and
every vertex of $R$ is $a$-critical by the global assumption.  This
contradicts \Cref{prop:25}.
\end{proof}

\section{Certificate semantics and coverage}\label{sec:certificate}

The finite computation is organised around a forced two-coloured pentagon.  In
a good two-colouring of $K_5$, both colour degrees are at most two and sum to
four, so both colours are pentagons.  At order $15$ or $16$, select a vertex
and five same-colour neighbours to obtain such a pentagon.  Fix its rim and
diagonal colours, name the third colour, and sort the remaining vertices by
their ordered edge colours into the pentagon.

\begin{computation}[Catalogue completeness and extension]\label[computation]{comp:catalogues}
The normalised catalogues contain $66$ labelled entries at order $15$ and six
at order $16$.  Every entry was checked to be a good colouring with its
prescribed labels.  Blocking each entry and replaying the two DRAT proofs shows
that the catalogues are complete under the chosen symmetry and sorting
normalisation.  Deleting each possible unnamed vertex of every order-$16$
entry produces exactly the $66$ order-$15$ entries.  Thus the extension property
of \Cref{lem:extend} holds.
\end{computation}

The counts $66$ and six are counts under this labelled normalisation, not
counts of isomorphism classes.  The sorting auxiliaries are one-way
prefix-equality indicators: when preceding coordinates agree they enforce
nondecreasing colours.  No false converse is imposed.  A physically sorted
colouring therefore admits the auxiliary assignment needed by the CNF.

For the 25-vertex obstruction, the common five-set $F=L\cap R$ omits both $a$
and $b$, so its other two colours again form rim and diagonal pentagons.  Use
one common order for $F$ on both arms.  Selecting one catalogue entry per arm
gives $66^2=4,356$ placements.  The simultaneous dihedral action on $F$ and
exchange of the two arms, including $a$ and $b$, reduce these to $247$ orbits.

\begin{computation}[Two complete refutations per orbit]\label[computation]{comp:orbits}
All $247$ orbits were refuted in both a slot encoding and a word encoding.
The verifier reconstructs the orbit decomposition and checks that these orbits
partition all $4,356$ placements.  Together with the two catalogue proofs, the
certificate therefore consists of
\[
 2+247+247=496
\]
distinct replayed DRAT refutations.  Its replay record reports 4,356
placements, 247 orbits and 496 checked proofs.
\end{computation}

The direction of the encoding implications is the essential semantic point.

\subsection{Slot encoding}

For a vertex $v$ in one arm, let $c$ be the colour omitted by the opposite arm
and let $Q$ be its colour-$c$ neighbours inside its own arm.  There are four or
five such anchors.  If $d_c(v)\geq15$, its entire colour-$c$ neighbourhood has
order at least $15$ and, by the catalogue, embeds into an order-$16$ template.
Every further physical vertex joined to $v$ in colour $c$ must occupy a
distinct compatible template position.

The slot encoding chooses a template and position for each such vertex,
prescribes the edge colours to the anchors, requires distinct positions and
enforces compatibility between pairs of entering vertices.  Given a genuine
colouring satisfying the hypotheses of \Cref{prop:25}, the actual embeddings
supply all auxiliary choices.  Thus UNSAT of this necessary-condition formula
excludes the ambient colouring.

\subsection{Word encoding}

The word encoding records vectors of edge colours into $Q$ together with the
compatible pair colours.  With four anchors, distinct positions can have the
same vector; the encoding therefore allows every pair colour obtainable from
distinct compatible positions.  It too admits the actual embedding as an
auxiliary assignment.  Among the 247 orbits, 230 already close using the
five-anchor tests, while the remaining 17 also use the four-anchor tests.

A satisfying assignment to either relaxation would not, by itself, construct a
good ambient colouring.  Conversely, a refutation suffices precisely because
every genuine colouring yields an auxiliary extension.  The two encoders were
implemented separately and both exclude every orbit.  They are redundant
cross-checks, not independent proofs of the whole theorem: both use the same
catalogues, orbit cover and global reduction.  The slot refutations alone
suffice for the finite local obstruction.

\begin{remark}[Trust boundary]\label[remark]{rem:trust}
The proof of \Cref{thm:common} combines ordinary mathematical reductions with
finite SAT certificates and their checker.  The 496-refutation package was
replayed against regenerated inputs; it was not inferred from Kissat or
CaDiCaL exit statuses \cite{Kissat,Wetzler2014}.  This certificate is not yet a
Lean kernel proof of \Cref{thm:common}.  In contrast, the separate signed-unit
cover theorem in \Cref{sec:formal} is kernel-checked in Lean.
\end{remark}

\section{Reciprocity and the two profiles}\label{sec:reciprocity}

First collect the elementary degree bounds on $K_{58}$.  Since $N_c(v)$ omits
$c$, it has at most $R(3,3,3)-1=16$ vertices.  The four degrees at $v$ sum to
$57$, hence
\begin{equation}\label{eq:degrees}
 9=57-3\cdot16\leq d_c(v)\leq16,\qquad
 \sum_{c=0}^{3}d_c(v)=57.
\end{equation}

\begin{lemma}[Colour connectivity]\label[lemma]{lem:connected}
In every good four-colouring of $K_{58}$, each colour graph is connected.
\end{lemma}
\begin{proof}
Fix $c$ and suppose $u,v$ lie in different components.  The sets $N_c(u)$ and
$N_c(v)$ are disjoint: a common vertex would join the components.  Each set is
$c$-independent by triangle-freeness, and no cross edge has colour $c$, again
by disconnectedness.  Their union is therefore $c$-independent of size at
least $2\cdot9=18$, contradicting the upper bound $16$.
\end{proof}

\begin{proof}[Proof of \Cref{cor:profiles}]
Under reciprocity, criticality in colour $c$ is constant along every $c$-edge.
By \Cref{lem:connected}, it is constant on all $58$ vertices, separately for
each $c$.  Every vertex has at least one critical colour, since otherwise its
degree sum would be at most $4\cdot14<57$.  \Cref{thm:common} allows at most
one common critical colour, so there is exactly one, and it is the same colour
at every vertex.

At any vertex, this colour has degree $15$ or $16$ by \eqref{eq:degrees}.  The
other three degrees are at most $14$.  If the common critical degree is $15$,
they sum to $42$, forcing $(14,14,14)$.  If it is $16$, they sum to $41$ and
are at most $14$, forcing $(14,14,13)$.  Sorting gives the two displayed
profiles.
\end{proof}

For comparison, \Cref{tab:profile-arithmetic} lists all eleven numerical
profiles, rather than only the seven profiles containing the current OPEN
wall.  The last column is the consequence \emph{if reciprocity is assumed};
it does not change the unrestricted ledger status.

\begin{table}[htbp]
\centering
\caption{All sorted degree profiles at $58$: the four degrees sum to $57$ and
are at most $16$.  A critical colour has degree at least $15$.}
\label{tab:profile-arithmetic}
\begin{tabular}{lcl}
\toprule
Profile & Number of critical colours & Under reciprocity\\
\midrule
$(16,16,16,9)$  & 3 & excluded\\
$(16,16,15,10)$ & 3 & excluded\\
$(16,16,14,11)$ & 2 & excluded\\
$(16,16,13,12)$ & 2 & excluded\\
$(16,15,15,11)$ & 3 & excluded\\
$(16,15,14,12)$ & 2 & excluded\\
$(16,15,13,13)$ & 2 & excluded\\
$(15,15,15,12)$ & 3 & excluded\\
$(15,15,14,13)$ & 2 & excluded\\
$(15,14,14,14)$ & 1 & not excluded by this argument\\
$(16,14,14,13)$ & 1 & not excluded by this argument\\
\bottomrule
\end{tabular}
\end{table}

\section{Proof of the ticket gap}\label{sec:ticket-gap}

Throughout this section the host is a good four-colouring of $K_{58}$.  Split
each colour-$c$ graph at the high/low cut $(H_c,X_c)$ defined in
\Cref{sec:consequences}.  A colour is \emph{mixed} when both sets are nonempty.
Only crossing edges carry unreciprocated tickets.

The first step strengthens the two-critical theorem when one colour has very
few low vertices.

\begin{lemma}[Low-set incidence bound]\label[lemma]{lem:incidence}
If $|X_a|\leq2$, then
\begin{equation}\label{eq:incidence}
16|X_b|\ \geq\ \sum_{r\in H_a}(d_a(r)-14)
   =|H_a|+|\{r:d_a(r)=16\}|
\end{equation}
for every colour $b\neq a$.  In particular, on $58$ vertices,
$|X_a|\leq2$ implies $|X_b|\geq4$.
\end{lemma}
\begin{proof}
We first prove that $|N_a(r)\cap H_b|\leq14$ for every $r$.  Otherwise choose
a 15-set $L$ in this intersection.  As before, $L$ omits $a$ and extends to a
good three-coloured $K_{16}$.  Let $Y\subseteq L$ be the ten vertices of
internal colour-$b$ degree five.  The catalogue computation checks that every
five-subset of $Y$ contains a $b$-edge.  Consequently, a vertex $x\in X_a$ can
be $b$-joined to at most four vertices of $Y$, since $N_b(x)$ is $b$
independent.  There are at most two such $x$, so some $y\in Y$ has no $b$-edge
to $X_a$.

Choose $15$ colour-$b$ neighbours of $y$ in $H_a$, including its five inside
$L$, to form $R$.  This is possible because at most two ambient $b$-neighbours
lie outside $H_a$ and none of the five internal neighbours does.  The sets
$L,R$ then violate \Cref{prop:25}: $L\subseteq H_b$ and $R\subseteq H_a$.

For each $r\in H_a$, at least $d_a(r)-14$ of its colour-$a$ neighbours lie in
$X_b$.  Summing over $r$ counts each colour-$a$ edge into $X_b$ at most once at
its $H_a$ endpoint and at most $16$ times at its $X_b$ endpoint.  This gives
\eqref{eq:incidence}.  Finally, $d_a\leq16$ makes the summands $1$ or $2$,
which proves the displayed equality, and $|H_a|\geq56$ gives the last claim.
\end{proof}

\begin{lemma}[Size of a mixed colour cut]\label[lemma]{lem:cut}
The number of colour-$c$ edges crossing $(H_c,X_c)$ is between $9$ and $14$
when $|X_c|=1$, and is at least $15$ whenever $|X_c|\geq2$.
\end{lemma}
\begin{proof}
The singleton assertion is just \eqref{eq:degrees}.  Suppose now that
$\ell=|X_c|\geq2$ and at most $14$ edges cross.  Call a vertex touched when it
is incident to a crossing edge.

There is an untouched high vertex.  Otherwise $h=|H_c|\leq14$.  Each high
vertex has at least $16-h$ crossing neighbours, so at least $h(16-h)\geq15$
edges leave $H_c$, a contradiction.  There is also an untouched low vertex.
Otherwise $2\leq\ell\leq14$.  Since the $c$-graph on $X_c$ is triangle-free,
it has at most $\lfloor\ell^2/4\rfloor$ internal edges.  Summing degrees gives
at least
\[
 9\ell-2\left\lfloor\frac{\ell^2}{4}\right\rfloor\geq16
\]
edges leaving $X_c$, again a contradiction.

Choose an untouched high $u$.  It has at least $15$ high neighbours and at
most $14$ touched neighbours, so choose an untouched high neighbour $v$.
Choose also an untouched low vertex $w$.  The independent sets $N_c(u)$ and
$N_c(w)$ lie on opposite sides and have at least $15+9=24$ vertices together.
At least eight edges join them: deleting one endpoint from each of at most
seven joining edges would leave an independent set of size at least $17$.
The same count applies to $N_c(v)$ and $N_c(w)$.  Since $uv$ has colour $c$,
triangle-freeness makes $N_c(u)$ and $N_c(v)$ disjoint.  These are at least
$16$ different crossing edges, contradicting the assumption of at most $14$.
\end{proof}

\begin{proof}[Proof of \Cref{cor:ticket-gap}]
If only one colour is critical anywhere, then \eqref{eq:degrees} forces it to
be critical everywhere, and no colour cut is mixed.  Hence all tickets are
reciprocated.  Suppose instead that at least two colours are critical
somewhere.  \Cref{thm:common} implies that at least one colour is mixed.

If two or more colours are mixed, \Cref{lem:cut} contributes at least
$2\cdot9=18$ unreciprocated tickets.  If exactly one colour is mixed, another
colour $a$ is critical everywhere.  Applying \Cref{lem:incidence} to $a$ shows
that the mixed colour has at least four low vertices, so its cut contributes
at least $15$.  This proves the first two assertions.

For the final assertion, suppose that incoming and outgoing ticket totals
balance at every vertex.  Reciprocated edges have no net directed
contribution.  If two or more colours are mixed, they cannot all have singleton
low sets: two singleton low sets in colours $a,b$ would contradict
\Cref{lem:incidence}.  Some mixed colour therefore has at least two low
vertices and sends at least $15$ net outgoing tickets across its high cut.
Vertexwise balance requires at least $15$ further unreciprocated tickets in
the opposite direction.  This proves the lower bound $30$.
\end{proof}

\section{The frozen cover at 58}\label{sec:cover}

We now turn to the larger campaign, which is separate from the general
obstruction above.  It follows the global ticket strategy of
\cite{Tatarevic2026}, adapted to $58$ vertices.  A vertex $u$ sends a ticket
to $v\neq u$ when $d_{c(uv)}(u)\geq15$.  If $H(v)$ is the number of tickets
received by $v$ and $w(v)$ its total number sent, then
\[
 w(v)=\sum_{c:d_c(v)\geq15}d_c(v),\qquad
 \sum_v H(v)=\sum_v w(v).
\]
Consequently some root satisfies $H(r)\geq w(r)$.

At a chosen root, a finite table of 41 threshold plans is used.  The selected
plan satisfies a slack inequality against $w(r)$.  If every arm stayed below
its threshold, $H(r)$ would be below $w(r)$.  Thus one arm reaches its
threshold and contains a sender, called the mark.  Classification of a critical
neighbourhood at the mark, palette normalisation and support embeddings turn
that configuration into signed edge facts.  The frozen enumeration retains
113,179 case keys after symmetry reduction.

The surrounding audit recomputed the arithmetic of the enumeration:
\[
 190,736\ \text{raw placements}
 =113,179\ \text{retained cases}+77,557\ \text{transported placements}.
\]
It covers 1,764 families and their 41 plan strata.  Every omitted placement
is transported by explicit vertex and palette permutations, rather than merely
by an appeal to symmetry.  A generated Lean replay checks this finite object;
\Cref{sec:formal} states its exact conclusion.

There are three logically distinct layers:
\begin{enumerate}[label=\textup{(\arabic*)},nosep]
 \item \textbf{Cover completeness:} every good colouring satisfies the signed
 units of at least one frozen case.
 \item \textbf{Encoding adequacy:} that colouring also satisfies every clause
 of the corresponding emitted CNF, after assigning auxiliary variables.
 \item \textbf{Computational exclusion:} every one of those CNFs is
 unsatisfiable, with checkable evidence.
\end{enumerate}
A complete proof of $\R\leq58$ needs all three.  The present project proves
(1), has proved only part of (2), and has only a partial dataset for (3).

\section{Lean formalisation and the encoding boundary}\label{sec:formal}

The formal development uses Lean 4.33.1 and Mathlib \cite{Lean2021,Mathlib2020}.
Its public axiom-audit list admits only subsets of
\texttt{propext}, \texttt{Classical.choice} and \texttt{Quot.sound}; proof
placeholders and \texttt{native\_decide} are rejected.  Standard finite
\texttt{decide} computations are kernel reductions, not additional trust
axioms.

At the mathematical core, theorems \path{two_colour_six},
\path{three_colour_seventeen} and \path{deg_le_sixteen} establish the
small Ramsey inputs and degree cap.  \path{root_plan_arm} implements the
balanced-root threshold step.  An earlier formulation,
\path{exists_arm_reaches}, was retained only as an explicit record of a
mistake: its universally quantified strict inequality was impossible at the
zero degree vector.  The replacement assumes the plan inequality at the actual
root profile and is nonvacuous.

The imported 15/16-point catalogues establish \texttt{classified\_15} and
\texttt{classified\_16}.  Their pinned upstream source is recorded by hash,
and their axiom dependencies are checked in the current Lean build.  This is
separate from the 66/6-entry certificate used for \Cref{thm:common}.

The endpoint
\[
\texttt{K58.Replay.frozenCover : K58.FrozenCover}
\]
states that every good four-colouring of $K_{58}$ realises at least one of the
113,179 frozen signed-unit cases.  The recorded generated replay contains
1,588 modules and checks 1,764 families, 190,736 raw placements and 77,557
transports.  Its endpoint reports only the three standard axioms above.  It is
not an UNSAT theorem.

Encoding adequacy remains the principal formal gap.  The deployed cardinality
constraints were compared clause for clause with a staggered sequential
counter, including variable allocation.  The Lean theorem
\texttt{satisfiable\_atLeast} proves that, if sufficiently many input atoms
are true, the counter auxiliaries can be assigned without changing those
atoms.  Guarded cardinality wrappers, support selectors, the assembly of all
clause families and the hash-level binding to emitted DIMACS still need to be
completed.  In particular, refuting an arbitrary relaxation or a file that
merely resembles a case is not enough.  The intended implication must be
\[
 \text{physical colouring}\Longrightarrow
 \text{full CNF has an auxiliary extension},
\]
so that UNSAT transfers back to the colouring.

None of the 496 certificates for \Cref{thm:common} has yet been converted into
a Lean proof.  The supplementary file
\path{paper/lean/ReciprocityBoundary.lean} nevertheless checks the conditional
profile argument in the kernel.  Its theorem
\path{reciprocal_profile_with_certificate} takes triangle-freeness, reciprocity
and an explicitly defined \path{CommonCriticalCertificate} as premises, proves
colour connectivity, and returns exactly the two profiles.  It neither imports
nor asserts the finite certificate as a theorem.  Conversely, the proof of
\Cref{thm:common} does not use the missing $K_{58}$ CNF bridge.

\section{Computational status on 25 November 2026}\label{sec:status}

\Cref{tab:ledger} reproduces the cross-tabulation from the latest release audit.
The first column is an owner-reported status: the release reports that 98,142
CNFs were regenerated and their DRAT certificates replayed with zero failures,
manifest
\path{98d82268ff4c277820f00cbed3405c625dbb53d3603610e45eacc73b562b19f9}.
The release itself contains indices and recipes, not those proof bytes.  The
local audit checked the manifest and case identities, matched the frozen
ledger byte for byte, and rebuilt 12 sampled certified CNFs to their recorded
hashes.  It could not independently replay certificates it did not receive.

\begin{table}[htbp]
\centering
\caption{Latest release-index snapshot.  ``Certified'' is the owner-reported
single-DRAT pipeline status; ``legacy'' groups historical evidence and scoped
exclusions.  These are unconditional ledger fields, deliberately unaffected by
\Cref{cor:profiles}.}
\label{tab:ledger}
\begin{tabular}{lrrr}
\toprule
Profile & Certified & OPEN & Legacy\\
\midrule
$(15,15,15,12)$ & 82,072 & 1,288 & 80\\
$(16,15,15,11)$ & 16,070 & 20 & 83\\
$(15,15,14,13)$ & 0 & 7,732 & 0\\
$(16,15,14,12)$ & 0 & 2,058 & 0\\
$(16,15,13,13)$ & 0 & 1,046 & 0\\
$(15,14,14,14)$ & 0 & 381 & 0\\
$(16,16,14,11)$ & 0 & 140 & 0\\
$(16,16,13,12)$ & 0 & 140 & 0\\
$(16,14,14,13)$ & 0 & 88 & 0\\
$(16,16,15,10)$ & 0 & 0 & 1,767\\
$(16,16,16,9)$ & 0 & 0 & 214\\
\midrule
Total & 98,142 & 12,893 & 2,144\\
\bottomrule
\end{tabular}
\end{table}

The OPEN total consists of 1,308 rows in the two large profiles and 11,585 in
seven small profiles.  Twenty of the 1,308 have separate split certificates
that the one-DRAT-per-case index does not represent.  The 2,144 legacy entries
comprise 1,743 historical rows whose 11.4\,GB proof store is not present and
401 scoped exclusions whose arguments need explicit evidence.  These categories
must not be silently folded into either ``proved'' or ``UNSAT unknown.''

The seven small profiles are not a routine retry tail.  Every sampled row in
them remained UNKNOWN at the budgets tried.  In the larger released batches
these probes total 0 closures from 98 rows at $12$ seconds, 68 at $70$ seconds,
51 at $60$ seconds and nine at $300$ seconds.  Such experiments measure solver
behaviour only.  They do not establish satisfiability or unsatisfiability, and
they are not used as premises in \Cref{thm:common,cor:profiles}.

Of the reported certificates, 76,967 use the ordinary recipe and 21,175 use
the \texttt{full\_root} recipe.  The latter adds the balanced-root hypothesis
$H(r)\geq w(r)$, which the campaign established from the case's charge pin.
That distinction is part of the evidence, not an implementation detail.

\section{Audit methodology and remaining obligations}\label{sec:audit}

The project used several independent checks of the supplied reduction rather
than assuming that a complete-looking bundle was sound.  The useful principles
were as follows.

\begin{enumerate}[label=\textup{(\alph*)},leftmargin=2em]
 \item \textbf{Regenerate and bind the object.} Case identifiers were
 recomputed from definitions, and sampled formulas were regenerated from the
 frozen row and recipe.  A hash of a filename is not a substitute for this.
 \item \textbf{Replay proofs, not solver claims.} A refutation counts only
 after an independent checker accepts a proof against the regenerated input.
 Solver timeouts and exhausted heuristic searches remain UNKNOWN.
 \item \textbf{Check necessity in the right direction.} For auxiliary
 encodings, a real colouring must determine an auxiliary extension.  Extra
 constraints that a real colouring need not meet invalidate an exclusion even
 if every generated CNF has a perfect proof.
 \item \textbf{Test vacuity and negative controls.} The impossible zero-vector
 hypothesis in the early threshold theorem is an example of an error that a
 successful compilation alone does not expose.  Boolean split coverage and
 modified-file negative controls played the same role in the computational
 checks.
 \item \textbf{Separate provenance levels.} Kernel theorems, locally replayed
 DRAT certificates, owner-reported replays, historical evidence references
 and heuristic measurements are recorded distinctly.
\end{enumerate}

The supplied-package audit found defects spanning disconnected theorem
citations, incomplete recipe fields, checker invocations that were not actually
performed, acceptance of a hand-written manifest without executing a checker,
optional support tables despite their use in completeness, and inadequate
per-case encoding pins.  These are findings about the supplied tooling, not
claims that every current script retains the same defects.  They explain why
this paper does not equate a green manifest with a proof.

The following obligations remain before any bound $\R\leq58$ can be stated:
\begin{enumerate}[label=\textup{(\arabic*)},nosep]
 \item finish and check the full encoding-adequacy and DIMACS-binding layer;
 \item resolve or explicitly discharge the 12,893 OPEN entries;
 \item supply the historical proof archives or independently refute their
 rows, and replace the 401 scoped exclusions with complete evidence;
 \item assemble the resulting finite computational premise and its semantic
 theorem without introducing an unexamined axiom.
\end{enumerate}
Independently, a pure-kernel formalisation of the 496-certificate proof of
\Cref{thm:common} would strengthen its trust boundary, but is not required for
the honest conditional statement in \Cref{cor:profiles}.

\subsection*{Source and reproducibility note}

The project sources and cited audit records are at
\url{https://github.com/jiji7879/r3333-58}.  The principal internal records
for this draft are
\path{results/common-critical-audit-2026-11-25/AUDIT.md},
\path{results/k58-full-cover-run-2026-11-25/README.md},
\path{review/LEAN_FROZEN_COVER_VERIFIED.md},
\path{review/LEAN_ENCODING_ADEQUACY.md}, and
\path{review/ERRATA.md}.  The manuscript's exact rerun commands and outcomes
are recorded separately in \path{paper/VERIFICATION.md}; its numerical sanity
checks are in \path{paper/check_claims.py}.  Stored scratch directories,
absolute machine paths and unattached proof archives are not treated as
permanent evidence.

\section*{Acknowledgements}

This working draft uses mathematical arguments, formalisations and audit
artifacts developed in the r3333-58 project, including a supplied
common-critical certificate package.  Human authorship, contribution and
responsibility statements must be supplied before external submission.

\bibliographystyle{plain}
\bibliography{references}
\end{document}
